\documentclass{article}
\usepackage[utf8]{inputenc}
\usepackage[italian]{babel}
\usepackage[ruled,vlined]{algorithm2e}
\usepackage{amsmath}

\begin{document}

\section*{Introduzione}

Il \textbf{BucketSort} nello specifico ordina dei \textit{record}, ovvero 
strutture contenenti più tipi di dati. La sua tecnica si basa sui così detti
\textbf{bucket}, ovvero celle all'interno di una \textbf{struttura ausiliaria} che ci 
aiutano a ordinare gli elementi della struttura originale.

\vspace{0.3cm}
\begin{algorithm}[H]
\caption{\textsc{BucketSort}($A, n, k$)}
\KwIn{Array $A$ di lunghez za $n$, numero $k$ di bucket}

\tcp{Dove $n$ è la lunghezza di $A$ e $k$ è il numero di bucket}

\BlankLine

\tcp{Istanzio l'array ausiliario $Y$ di dimensione $k$}

\For{$i = 1$ \KwTo $k$}{
    $Y[i] \gets \text{coda vuota}$
}

\BlankLine

\For{$i = 1$ \KwTo $n$}{
    \If{$A[i].\text{chiave} \in [1 \dots k]$}{
        $Y[A[i].\text{chiave}].\text{enqueue}(A[i])$
    }
}

\BlankLine

\tcp{Contatore di $A$}
$j \gets 1$\;

\For{$i = 1$ \KwTo $k$}{
    \While{\textnormal{\textbf{not}} $Y[i].\text{empty}$}{
        $A[j] \gets Y[i].\text{dequeue}()$\;
        $j \gets j + 1$\;
    }
}
\end{algorithm}
\vspace{0.3cm}

L'algoritmo quindi popolerà la \textit{struttura ausiliaria} $Y$ con elementi aventi la stessa radice 
oer poi ricopiarli, in maniera ordinata, nella \textit{struttura originale}.
Inoltre è una tecnica \textbf{stabile} poichè se due elementi con la stessa chiave finiscono
nello stesso bucket "usciranno" nello stesso ordine in cui sono entrati.

\section*{Costic computazionali}
L'algoritmo dipende soltanto da due variabili, $n$ e $k$, quindi il suo costo è semplicemente 
$O(n + k)$, ma se $k$ è direttamente proporzionale a $n$ allora i costo diventerà \textbf{lineare}:
$O(n)$, \\
Allo stesso modo però se ci ritroviamo con un $k = n^2$ il costo si alzerà ancora di più. 


\end{document}